\documentclass[12pt]{article}
\usepackage{amsmath,amssymb}
\newcommand{\Dr}{\operatorname{Dr}}\newcommand{\Trop}{\operatorname{Trop}}\newcommand{\rk}{\operatorname{rk}}\newcommand{\cl}{\operatorname{cl}}
\begin{document}
\section*{Research notes}

Problem: Given a matroid $M$ and the relative Dressian $\Dr(M)$ of tropical
linear subspaces contained in $\Trop(M)$, suppose the lattice of flats
$\mathcal L(M)$ has an order-reversing involution $F\mapsto F^\perp$. Need
an order-reversing involution on $\Dr(M)$ extending
$F\mapsto F^\perp$ under the embedding
$F\mapsto \Trop(M/F\oplus U_{0,F})$.

Key structural point: a finite geometric lattice with an anti-automorphism is
modular. Geometric lattices are upper semimodular; applying the anti-automorphism
gives lower semimodularity, hence rank equality
\[
r(X)+r(Y)=r(X\wedge Y)+r(X\vee Y)
\]
for all $X,Y$. In a semimodular lattice this means every pair is modular, so the
lattice is modular. Thus the hypothesis is much stronger than mere
self-duality of the matroid.

Valuated quotient model: a tropical linear subspace
$\Trop(\nu)\subseteq \Trop(M)$ is a valuated quotient
$\nu\leftarrow \mu_M$, where $\mu_M$ is the trivial valuation of $M$.
Inclusion of tropical linear spaces corresponds to the quotient order.

Important lemma used in final proof: if $\nu$ is a rank $k$ valuated quotient of
$M$, then its value is constant on bases spanning the same rank-$k$ flat of $M$.
Proof idea: for adjacent bases $S+b$ and $S+b'$ of the same flat $X$, choose a
basis $C$ of $M/X$ and apply the incidence-Plucker relation indexed by
$(S,S\cup\{b,b'\}\cup C)$. Only the two terms $b,b'$ are finite, forcing
equality.

Construction: for a rank $k$ quotient $\nu$ define a rank $r-k$ quotient
\[
\nu^\perp_M(A)=
\begin{cases}
\nu(B),& A\text{ independent and } \cl(B)=(\cl A)^\perp,\\
\infty,&\text{otherwise}.
\end{cases}
\]
Well-defined by the previous lemma. The modular flat lattice makes the
Plucker and incidence-Plucker relations invariant under opposition
$X\mapsto X^\perp$, so $\nu^\perp_M$ is again a valuated quotient.
For a flag $\theta\leftarrow\nu\leftarrow\mu_M$, the same relation check gives
$\nu^\perp_M\leftarrow\theta^\perp_M\leftarrow\mu_M$, hence order reversal.
Involutivity is immediate from $(X^\perp)^\perp=X$.

Flat check: for $q_F=M/F\oplus U_{0,F}$ and $Y=\cl(A)$ with
$\rk(Y)=\rk(F)$, $\Phi(q_F)(A)$ is finite iff $Y^\perp\vee F=\hat1$.
By modularity and rank complement this is equivalent to
$Y^\perp\wedge F=\hat0$, equivalently $Y\vee F^\perp=\hat1$, the basis
condition for $M/F^\perp\oplus U_{0,F^\perp}$.

Potential examples considered:
- Boolean/free matroid: recovers usual valuated matroid duality
$p^*(A)=p(E\setminus A)$, possibly composed with a coordinate permutation.
- Rank-3 projective planes/near-pencils: points in $\Trop(M)$ are valuations on
points satisfying line minima; codimension-one subspaces are valuations on
lines satisfying dual pencil minima. A polarity sends one to the other.
- Nonmodular self-dual matroids such as V\'amos are irrelevant because a flat
lattice anti-automorphism would force modularity; matroid duality is not a flat
lattice anti-automorphism in general.

Citations useful:
Brandt--Eur--Zhang, Tropical flag varieties, Adv. Math. 384 (2021), for the
valuated flag matroid/quotient and tropical linear space inclusion
correspondence.
Dress--Wenzel, Valuated matroids, Adv. Math. 93 (1992).
\end{document}
